\section{Amenazas a la validez}

\textbf{Validez de constructo.} La métrica principal (exactitud estricta)
exige coincidencia textual con el ground truth bajo un protocolo que declara
explícitamente los sinónimos de valores categóricos cerrados como violación de
formato (Sección \ref{sec:prompting}). Esta amenaza ya no es sólo una
advertencia sin cuantificar, como en la ronda inicial de este experimento: queda
\textbf{resuelta parcialmente} por la exactitud laxa y por la taxonomía de
errores de la Sección \ref{sec:analisis-errores}, que separan explícitamente
cuánta exactitud estricta perdida corresponde a \texttt{uso\_de\_sinonimos} o
a \texttt{sin\_error\_semantico} en lugar de a un error de contenido real. La
brecha sigue siendo una cota superior semántica, no una medición
independiente con un evaluador humano, y la Sección 4.4 mostró que esta
métrica laxa es a su vez poco confiable donde la validez de formato es débil:
24 de las 27 respuestas que la métrica laxa acreditó a
\texttt{SmolLM2-360M-Instruct} tenían JSON inválido (ver más abajo, amenaza de
validez interna sobre el sesgo del juez y la fragilidad de
\texttt{SmolLM2-360M-Instruct}).

\textbf{Validez interna.}

\emph{Sesgo de auto-favorecimiento del juez.} El juez de la etapa 2 es
\texttt{granite-4.0-1b}, seleccionado por la regla congelada de mayor
exactitud estricta en la etapa 1 (Sección \ref{sec:dos-etapas}), y ese mismo
modelo es uno de los doce evaluados: \textbf{juzga también sus propias
respuestas de la etapa 1}. Un juez basado en LLM puede favorecer respuestas
generadas por sí mismo o por modelos de su misma familia
\cite{zheng2023judging}, y en este diseño esa posibilidad no está controlada
ni por un juez externo ni por un segundo evaluador independiente. El riesgo es
mayor todavía porque, al ser el modelo con mayor exactitud estricta de la
etapa 1, \texttt{granite-4.0-1b} es también el modelo con \textbf{menos}
respuestas incorrectas en juego (apenas 3 de 32, Tabla \ref{tab:taxonomia}):
tiene la menor cantidad relativa de casos propios para favorecer, lo que
acota — sin eliminar — el margen concreto de ese sesgo sobre el resultado
agregado. Además, es un modelo de menos de 2 mil millones de parámetros
adjudicando equivalencia semántica en lenguaje natural, una tarea de juicio
más exigente que la clasificación de campos estructurados, con la capacidad
limitada que ese tamaño implica; los 0 casos de \texttt{juez\_parse\_ok = False}
reportados en la Sección \ref{sec:analisis-errores} indican que el juez sí
logró producir salidas parseables de forma consistente, pero no dicen nada
sobre la calidad de su criterio semántico. Mitigar esta amenaza en trabajo
futuro requeriría un juez externo de mayor escala, no perteneciente al roster
evaluado, o una validación humana sobre una muestra de sus adjudicaciones.
Como contrapeso, la adjudicación es \textbf{reproducible por diseño}: ambas
etapas corren con decodificación greedy (\texttt{do\_sample=False},
temperatura 0), de modo que, dado el mismo modelo de entrada, cualquiera puede
re-ejecutar el juicio y obtener exactamente el mismo resultado.

\emph{Versiones divergentes de la librería de inferencia como confusor de
latencia.} Los doce modelos del roster activo corren bajo \textbf{dos grupos
de versión} de \texttt{transformers}, mutuamente excluyentes (Tabla
\ref{tab:versiones}): el grupo A resuelve a la versión 4.57.6, con
\textbf{9} modelos, y el grupo B resuelve a la versión 5.14.1, con los
\textbf{3} restantes. Esto es un confusor conocido para cualquier comparación
de \textbf{latencia} entre un modelo del grupo A y uno del grupo B: parte de
la diferencia observada puede deberse a la versión de la librería y no al
modelo en sí. Cada grupo está forzado por evidencia empírica propia, no por
preferencia de diseño: \texttt{granite-4.0-350m} necesita el grupo A porque
falla bajo 5.14.1; \texttt{LFM2.5-230M}, \texttt{LFM2.5-350M} y
\texttt{Qwen3.5-0.8B} necesitan el grupo B porque fallan bajo 4.57.6. La
evidencia más nítida de que el confusor no es un artefacto teórico es que el
grupo B resuelve exactamente a la versión (5.14.1) que rompe al modelo que
fuerza el grupo A: ninguna versión mayor única cubre el roster completo. No
se repite aquí la premisa, ya descartada en este mismo trabajo, de que una
única versión mayor de \texttt{transformers} elimina este confusor: la
evidencia de ambos grupos activos la refuta en las dos direcciones.

\emph{Topología de procesadores no observable.} La Sección
\ref{sec:entorno-repro} estableció que la configuración de Docker de esta
máquina expone únicamente 2 procesadores lógicos a los contenedores, y que
\texttt{-{}-cpuset-cpus=0-1} fija ese mismo par en las doce corridas y en la
del juez. Qué relación guardan esos 2 procesadores lógicos con los núcleos
físicos del host —hermanos de un mismo núcleo por \emph{hyperthreading}
(SMT), lo que introduciría contención adicional entre ellos, o dos núcleos
físicos distintos, sin esa contención— la decide el hipervisor de esa
configuración de Docker, y esta amenaza a la validez no puede resolverse
observando desde adentro del contenedor: la topología efectiva no quedó bajo
control de este experimento. Esto afecta por igual a los doce modelos —todos
corrieron bajo el mismo tope, de a uno por vez—, así que no introduce un
sesgo relativo entre ellos, pero sí implica que los valores absolutos de
latencia reportados pueden estar inflados por contención respecto de lo que
esta misma máquina lograría con dos núcleos físicos completos, y no deben
leerse como la mejor latencia alcanzable en este hardware. Los cocientes de
latencia entre variantes híbridas y densas de la Sección
\ref{sec:discusion-latencia-hibrida} (\texttt{$\times$4,10} y
\texttt{$\times$2,88}) son robustos a esta amenaza precisamente porque
comparan pares de modelos medidos bajo la misma topología, cualquiera haya
sido.

\textbf{Validez externa.}

Se conservan además las amenazas de validez externa de la ronda inicial de
este experimento: el dataset de evaluación es chico y desbalanceado (32
comandos, entre 3 y 8 por categoría lingüística), se usó un único prompt sin
optimización específica por modelo, la entrada fue siempre texto limpio sin
ruido de reconocimiento de voz (ASR), y sólo se cubrió una variante
dialectal (español rioplatense).

\textbf{Validez de conclusión.} Con entre 3 y 8 comandos por categoría, las
diferencias porcentuales entre modelos de la Tabla \ref{tab:categorias} y de
la Tabla \ref{tab:globales} son indicativas, no estadísticamente
significativas en sentido formal — la misma limitación de la ronda inicial
de este experimento, que este trabajo no resuelve. Todas las corridas del
barrido principal (Sección \ref{sec:entorno-repro}) se ejecutaron en una
\textbf{única máquina}, de forma estrictamente secuencial y con
procesadores fijos, para que la latencia fuera comparable entre modelos;
las corridas de control, en cambio, compartieron la máquina con trabajo
concurrente del orquestador y por eso se excluyen explícitamente de
cualquier afirmación de latencia de este trabajo, quedando limitadas a la
justificación metodológica de la Sección \ref{sec:entorno-repro}. Por último,
\texttt{attn\_implementation} y la \texttt{revision} exacta de cada modelo de
HuggingFace quedaron, también en este trabajo, \textbf{sin pinnear}: cada
corrida usó lo que \texttt{transformers} resolviera por defecto en el momento
de ejecutarse, y ningún registro deja constancia de qué implementación de
atención ni qué revisión de pesos se usó exactamente. Pinnear ambos valores
ahora exigiría re-correr el barrido completo para que la comparación siguiera
siendo de igual a igual, lo que explícitamente no se consideró que valiera la
pena frente al beneficio; se documenta aquí como limitación conocida y no como
un defecto corregido.
